\section{Introduction}\label{sec:introduction}

An insertion or deletion changes the positions of every later symbol in
a word. A long interval can therefore move at small edit cost even when
many positions change. This tension between order and inexpensive shifts
is central to the question studied here: how faithfully can one represent
the edit distances of an entire set of words by distances in $\ellone$?

We use ordinary edit distance $\ED(x,y)$: the minimum number of
single-symbol insertions, deletions and substitutions taking $x$ to $y$,
each with cost one. For a finite alphabet $\Sigma$ with at least two
symbols and an integer
$d\geq1$, let $\Sigma^{\leq d}$ be the set of words of length at most
$d$, including the empty word. The cap constrains the words in the
domain; an edit script computing their distance may pass through words
of any finite length.

For a finite metric space $(X,\rho)$ with at least two points and an
injective map $f:X\to\ellone$, where $\ellone$ is the space of absolutely
summable real sequences, define
\[
 \dist(f)=
 \left(\max_{x\ne y}\frac{\|f(x)-f(y)\|_1}{\rho(x,y)}\right)
 \left(\max_{x\ne y}\frac{\rho(x,y)}{\|f(x)-f(y)\|_1}\right).
\]
Let $c_1(X,\rho)$ be the infimum of $\dist(f)$ over such maps, and put
$E_\Sigma(d)=c_1(\Sigma^{\leq d},\ED)$. No restriction is imposed on
the dimension of the image or on the computation of the map. Throughout
the paper, logarithms in asymptotic estimates are natural unless a base
is written explicitly.

We prove the following lower bound in two independent ways. Each proof
uses a tree to combine inexpensive shifts with a displacement that no
alignment can hide.

\begin{theorem}[Binary lower bound]\label{thm:binary}
There are an absolute constant $c>0$ and an integer $d_0$ such that,
for every integer $d\geq d_0$, there exist an integer $1\leq n\leq d$ and a
finite set $\mathcal W\subseteq\{0,1\}^n$ with $|\mathcal W|\geq2$
satisfying
\[
 c_1(\mathcal W,\ED)
 \geq \exp\!\left(c\sqrt{\log d\,\log\log d}\right).
\]
\end{theorem}

The words in a witness all have the same length, although that length
may depend on the cap $d$. At tree depth $k$, each construction
produces logarithmic distortion of order at least $k\log k$ and
logarithmic word length of order $k^2\log k$. Balancing these two scales
gives Theorem~\ref{thm:binary} for every sufficiently large cap, not just
for a subsequence of caps.

The lower bound is proved here without a companion input. The uniform
upper embedding theorem of the finite-circle companion, stated precisely
in Section~\ref{sec:consequences}, then gives
\[
 \log E_\Sigma(d)=\Theta\!\left(\sqrt{\log d\,\log\log d}\right)
\]
with constants independent of the finite alphabet $\Sigma$. This is a
matching exponential scale, not a constant-factor estimate for
$E_\Sigma(d)$ itself.

\subsection{History and the embedding question}

The single-symbol error model is classical in coding and sequence
comparison. Levenshtein studied binary codes correcting insertions and
deletions \cite{Levenshtein1966}, and Wagner and Fischer formulated
string correction through increasing partial alignments
\cite[Sections~3 and~5]{WagnerFischer1974}. An embedding asks for one set of
coordinates that represents every pair in a domain. It is a different
task from estimating the distance of one given pair of words.
For supplied pairs of binary words, \cite[Theorem~1.1]{OpenAIEditApprox2026}
gives a randomized approximation scheme with almost-linear worst-case
expected time for each fixed rational accuracy.

Ostrovsky and Rabani proved an $\ellone$ distortion bound
$\exp(O(\sqrt{\log n\,\log\log n}))$ for binary words of a fixed length
$n$ \cite{OstrovskyRabani2007}. Their recursion compares overlapping
substrings and represents their distributions in smaller metric spaces.
They also noted extensions to larger alphabets and varying lengths.
For the upper bound used at the end of this paper, the companion uniform
upper embedding theorem
\cite[Uniform upper embedding theorem]{OpenAIUniformUpper2026}
gives the full quantitative statement on
$\Sigma^{\leq d}$ with constants independent of the finite alphabet.

On the lower-bound side, Andoni, Deza, Gupta, Indyk and Raskhodnikova
realized complete bipartite graph metrics in binary edit distance and
obtained $\ellone$ distortion approaching $3/2$
\cite[Section~2 and Theorem~3.1]{AndoniDezaGuptaIndykRaskhodnikova2003}.
Their witness uses three adjacent word lengths.
Khot and Naor developed Fourier methods for
nonembeddability and obtained an edit-distance lower bound of
$(\log n)^{1/2-o(1)}$ for binary words of a fixed length
\cite{KhotNaor2006}. Their insertion--deletion
metric and the ordinary metric used here differ by at most a factor of
two. Krauthgamer and Rabani then proved an $\Omega(\log n)$ lower bound
for binary words of length $n$ \cite[Theorem~1.3]{KrauthgamerRabani2009}.
Khot and Naor used Boolean noise sensitivity, while Krauthgamer and
Rabani sharpened the obstruction through an influence inequality for
bit flips and cyclic shifts. These cut and shift arguments are
methodological predecessors of the proofs here. A cut records which
pairs a binary coordinate separates, and Fourier coefficients describe
how that separation changes under translations. We prove the particular
translation estimates needed for our trees.

Alphabet reduction also has a substantial history. Schulman and
Zuckerman use separation between long intervals of distinct codewords
in concatenated insertion--deletion codes \cite{SchulmanZuckerman1999},
and synchronization strings protect positions against insertion and
deletion errors \cite{HaeuplerShahrasbi2017}. More directly,
Bhattacharya, Dey, Goldenberg and Kouck\'y state a map on all finite
words with an insertion--deletion guarantee normalized by the sum of
the word lengths, for pairs of positive total length. The block
substitution used to prove their stated theorem has logarithmic width
and, for a fixed error parameter, a constant target alphabet
\cite[Theorem~3; full version, Section~3.2.1 and Lemma~3.3]{BhattacharyaDeyGoldenbergKoucky2024}.
For equal-length inputs, their stated normalized guarantee, combined
with the elementary delimiter construction below, gives a binary
reduction with logarithmic width and absolute distortion. Our lower
proofs do not use that result. They use the short delimiter code, whose
width-dependent loss is harmless at the tree scale; the appendix gives
a separate constant-distortion code for one prescribed word length.
For comparison, the isometric Hamming-to-edit embedding of Bhattacharya
et al.\ \cite[Theorem~7]{BhattacharyaEtAl2026} starts from a different
input metric. The conversions proved here compare edit distance on the
source words with edit distance on their binary encodings.

\subsection{Two independent tree proofs}

Both constructions use an ordered tree. A leaf has a state coordinate
and produces a symbol carrying its leaf label. An internal node produces
a long list of rows, each containing its child words in their fixed
order. Advancing the row index deletes a few rows at one end and inserts
a few at the other, so it is inexpensive. A distinguished second shift
must instead separate the words even when an alignment crosses row
boundaries.

The shared alignment lemma in Section~\ref{sec:preliminaries} handles
this difficulty. Leaf labels preserve child indices, not row indices.
An unsplit source row loses symbols because its child words are
separated from the corresponding target words. A split source row
forces a long unmatched target interval between consecutive matches.
The intervals charged by different source rows are disjoint. Thus
separation in every comparison of two rows implies separation of the
whole words, without assuming that row boundaries match.

The two constructions arrange the shifts differently.
In the finite tree of Section~\ref{sec:finite}, a child's period is its
parent's period times a new prime. At an internal node, a vector $u_v$
advances rows and a multiple $w_v$ has the parent period. The distinct
prime factors prevent a difference of row indices from cancelling the
distinguished shift in too many children. This gives a deterministic
separation proof on a finite group. For a Fourier character, activity at
a node with few active children has a short period and is detected by
an average of inexpensive row shifts. A character escaping every such
test must be active on many leaves; individual leaf changes then bound
its contribution. The section closes with binary witnesses for every
large cap.

The circle tree in Sections~\ref{sec:torus-construction}
and~\ref{sec:torus-fourier} uses quantized phases instead of finite
residues. Random prime lists and coefficients make the root half-turn
separate every base phase. The probabilistic assertion is only for one
fixed shift; the all-phase conclusion follows deterministically from
the alignment lemma. In its Fourier argument, tiny rotations detect
large integer frequency sums and rational row shifts detect failed
prime congruences. When both tests fail, divisibility forces every
nonzero sum to continue through many children. This supplies the second
binary proof without using the finite tree.

Both routes use the same elementary delimiter code from
Section~\ref{sec:preliminaries}. Its distortion loss grows with the
block width, but that width is only $O(k\log k)$ here and does not
change the logarithmic lower-bound scale. Appendix~\ref{sec:code}
proves a stronger, separate alphabet-conversion theorem: for words of
one prescribed length $M$ over at most $M$ letters, blocks of width
$M^2$ preserve edit distance up to scale and an absolute factor $804$.
Its proof is independent of both trees, and neither headline lower
bound requires this stronger conversion.
